
\chapter*{Expos\'e \texorpdfstring{VII\textsubscript{B}}{VII-B}\\
Infinitesimal Study of Group Schemes}
\phantomsection
\label{sga3:VIIB:title}
\addcontentsline{toc}{chapter}{Expos\'e VII-B. Infinitesimal Study of Group Schemes}
\markboth{EXPOS\'E VII\textsubscript{B}. INFINITESIMAL STUDY OF GROUP SCHEMES}%
  {EXPOS\'E VII\textsubscript{B}. INFINITESIMAL STUDY OF GROUP SCHEMES}

\begin{center}
  \textit{by P. Gabriel}
\end{center}

\setcounter{footnote}{-1}
\section*{B) Formal Groups}
\phantomsection
\label{sga3:VIIB:B}
\addcontentsline{toc}{section}{B) Formal Groups}

\noindent\footnote{Editors' note: Version of 13 October 2024.}
The study of formal groups is usually extremely simple. If this does not
appear clearly in the pages that follow, the responsibility lies with an
arithmetician who claims to know formal groups over ``something other than
fields.''\footnote{The interest of formal groups over a complete
Noetherian local ring appears, for example, in the work of Lubin and Tate
(cf.\ \textup{[LT65]}). The study of formal groups over an arbitrary base,
and of lifting and deformation questions---in particular for
Barsotti--Tate groups (``\(p\)-divisible groups'')---has given rise to an
abundant literature; see, for example, \textup{[LT66, Ta67, Gr74, Me72,
La75, Fo77, Il85, Br00]}. In particular, the results of the present
expos\'e are largely taken up again in Chapter~I of \textup{[Fo77]}.}
We have therefore developed, for formal groups ``locally free over
projective limits of Artinian rings,'' the generalities ordinarily stated
for formal groups defined over a field. For a more detailed study of the
latter, we refer to the 1964/65 Heidelberg--Strasbourg algebraic geometry
seminar.\footnote{The editors found only such a seminar dated 1965/66 and
entitled ``Linear Algebraic Groups,'' in which the notion of a formal group
does not appear; see instead \textup{[De72]}.}

\section*{0. Review of Pseudocompact Rings and Modules}
\phantomsection
\label{sga3:VIIB:0}
\addcontentsline{toc}{section}{0. Review of Pseudocompact Rings and Modules}

This section contains a few technical preliminaries. We recall and
complete here several results from \textup{[CA]} (\emph{Des cat\'egories
ab\'eliennes}, Bull.\ Soc.\ Math.\ France \textbf{90} (1962)).

\medskip
\noindent\textbf{0.1.}
\phantomsection
\label{sga3:VIIB:0.1}
A \emph{left pseudocompact ring} is a separated, complete topological
ring with identity which has a basis of neighborhoods of \(0\) consisting
of left ideals \(\mathfrak l\) of finite colength, that is,
\[
  \operatorname{length}_A(A/\mathfrak l)<+\infty.
\]
We shall assume here that \(A\) is commutative, so that no distinction
need be made ``between left and right.''

In particular, the quotients \(A/\mathfrak l\) are Artinian rings, and
\(A\) identifies with the topological projective limit of these rings,
each endowed with the discrete topology.

A complete Noetherian local ring \((A,\mathfrak m)\) is plainly
pseudocompact.\footnote{When it is endowed with the
\(\mathfrak m\)-adic topology.}

\medskip
\noindent\textbf{0.1.1.}
\phantomsection
\label{sga3:VIIB:0.1.1}
Every closed ideal \(I\) of \(A\) is the intersection of the open ideals
that contain it.\footnote{Indeed, if \(x\notin I\), there is an open ideal
\(\mathfrak l\) such that
\((x+\mathfrak l)\cap I=\varnothing\); then \(I+\mathfrak l\) is an open
ideal that does not contain \(x\). In what follows, the editors have also
made explicit the fact that every ``closed maximal'' ideal is maximal and
open.}
Every closed maximal ideal is therefore open. Moreover, if
\(\mathfrak l\) is an open ideal of \(A\), the maximal ideals of
\(A/\mathfrak l\) correspond bijectively to the maximal ideals
\(\mathfrak m\) of \(A\) containing \(\mathfrak l\); the latter are thus
open and closed. Consequently, every closed maximal ideal is an open
maximal ideal (and hence closed), the converse being clear. We denote the
set of these ideals by \(\Upsilon(A)\).

If \(\mathfrak l\) is an open ideal of \(A\) and
\(\mathfrak m\in\Upsilon(A)\), the localization
\((A/\mathfrak l)_{\mathfrak m}\) is a local ring when
\(\mathfrak m\) contains \(\mathfrak l\), and is zero otherwise. Since
\(A/\mathfrak l\) is Artinian, it is a direct product of finitely many
local rings; thus
\[
  A/\mathfrak l
  \simeq
  \prod_{\mathfrak m\in\Upsilon(A)}
    (A/\mathfrak l)_{\mathfrak m}.
\]
It follows that there are ``canonical'' isomorphisms
\[
  A
  \simeq \varprojlim_{\mathfrak l}(A/\mathfrak l)
  \simeq \varprojlim_{\mathfrak l}
    \prod_{\mathfrak m}(A/\mathfrak l)_{\mathfrak m}
  \simeq \prod_{\mathfrak m}
    \varprojlim_{\mathfrak l}(A/\mathfrak l)_{\mathfrak m}
  \simeq \prod_{\mathfrak m}A_{\mathfrak m},
\]
where
\[
  A_{\mathfrak m}
  =
  \varprojlim_{\mathfrak l}(A/\mathfrak l)_{\mathfrak m}.
\]
This \emph{local component} \(A_{\mathfrak m}\) is a filtered
projective limit of Artinian local rings endowed with the discrete
topology. It is therefore a local ring which is pseudocompact for the
projective-limit topology.\footnote{Notice that \(A_{\mathfrak m}\) is
the localization of \(A\) at the maximal ideal \(\mathfrak m\). Indeed,
the identity element \(e_{\mathfrak m}\) of \(A_{\mathfrak m}\) is an
idempotent of \(A\) not belonging to \(\mathfrak m\). Since
\(A=A_{\mathfrak m}\times B\), where
\(B=A(1-e_{\mathfrak m})\), the ring \(A_{\mathfrak m}\) identifies with
the localization \(A_{e_{\mathfrak m}}\), and hence also with
\(S^{-1}A\), where \(S=A\setminus\mathfrak m\). Furthermore,
\(e_{\mathfrak m}(1-e_{\mathfrak m})=0\), so \(\mathfrak m\) contains
\(1-e_{\mathfrak m}\), hence also \(B\); thus \(\mathfrak m\) identifies
with \(\mathfrak mA_{\mathfrak m}\times B\).}

\medskip
\noindent\textbf{0.1.2.}
\phantomsection
\label{sga3:VIIB:0.1.2}
Let \(r(A)\) be the intersection of the open maximal ideals of \(A\),
that is, the cartesian product of the ideals
\(\mathfrak mA_{\mathfrak m}\) under the identification
\[
  A=\prod_{\mathfrak m}A_{\mathfrak m}.
\]
For every open ideal \(\mathfrak l\) of \(A\), the image of \(r(A)\) in
\(A/\mathfrak l\) is contained in the radical of \(A/\mathfrak l\).
Some power of this image is therefore zero, so
\(r(A)^n\subset\mathfrak l\) for all sufficiently large \(n\). Thus the
sequence \(r(A)^n\) tends to \(0\).

The same is true of \(x^n\) when \(x\in r(A)\). In other words, every
element of \(r(A)\) is topologically nilpotent, and the converse is
clear. It follows that the sequence with general term
\[
  1+x+\cdots+x^n
\]
converges, and converges to \(1/(1-x)\), whenever \(x\in r(A)\). This
shows that \(r(A)\) is the Jacobson radical of \(A\), that is, the
intersection of all maximal ideals of \(A\) (cf.\ Bourbaki,
\emph{Alg\`ebre}, Chap.~8, \S6, Th.~1).\footnote{Indeed, let
\(x\in r(A)\). If \(\mathfrak m\) were a maximal ideal not containing
\(x\), there would be \(y\in A\) such that
\(1-xy\in\mathfrak m\). But \(xy\in r(A)\), so \(1-xy\) is invertible,
a contradiction. Notice the following consequence: if
\(\Upsilon(A)=\{\mathfrak m_1,\ldots,\mathfrak m_r\}\) is finite, then
the \(\mathfrak m_i\) are all the maximal ideals of \(A\).}

\medskip
\noindent\textbf{Remarks.}\footnote{The editors added these remarks in
order to compare the definition of the formal spectrum
\(\operatorname{Spf}(A)\) given in 1.1 with those in EGA~I, 10.1.2 and
10.4.2.}

\noindent\textit{a)} If \(\mathfrak p\) is an open prime ideal of \(A\),
then, since \(A/\mathfrak p\) is Artinian, \(\mathfrak p\) is maximal.
Consequently, \(\Upsilon(A)\) is the set of open prime ideals of \(A\).

\smallskip
\noindent\textit{b)} Each \(\mathfrak mA_{\mathfrak m}\) is an ideal of
definition of \(A_{\mathfrak m}\), that is, an open ideal \(I\) for
which \(I^n\) tends to \(0\) (cf.\ EGA~0\textsubscript{I}, 7.1.2).
Consequently,
\[
  \operatorname{Spec}
    (A_{\mathfrak m}/\mathfrak mA_{\mathfrak m}),
\]
endowed with the topological ring \(A_{\mathfrak m}\), is an affine
formal scheme in the sense of EGA~I, 10.1.2.

\smallskip
\noindent\textit{c)} The topological ring \(A\) is admissible in the
sense of EGA~0\textsubscript{I}, 7.1.2 if and only if \(r(A)\) is open
(and hence an ideal of definition), and this holds if and only if
\(\Upsilon(A)\) is finite. In that case the affine formal scheme
\[
  \operatorname{Spf}(A)
  =
  \operatorname{Spec}(A/r(A))
\]
(cf.\ EGA~I, 10.1.2) has \(\Upsilon(A)\), endowed with the discrete
topology, as its underlying space; the sections of its structure sheaf
over a subset \(E\subset\Upsilon(A)\) form the ring
\[
  \prod_{\mathfrak m\in E}A_{\mathfrak m}.
\]

\smallskip
\noindent\textit{d)} Let \(A\) be an arbitrary pseudocompact ring. In
\SGARefLink{sga3:VIIB:1.1}{1.1} below, the space \(\Upsilon(A)\) is endowed with the discrete
topology and with the sheaf of rings whose ring of sections over any
subset \(E\) is
\[
  \prod_{\mathfrak m\in E}A_{\mathfrak m}.
\]
By \textit{b)}, every point then has an open neighborhood which is an
affine formal scheme; this therefore defines a formal scheme (EGA~I,
10.4.2), denoted by \(\operatorname{Spf}(A)\). For this formal scheme to
be affine, it must be quasi-compact, so \(\Upsilon(A)\) must be finite;
in that case \(\operatorname{Spf}(A)\) agrees with the definition of
EGA~I, 10.1.2.

\medskip
\noindent\textbf{0.1.3.}
\phantomsection
\label{sga3:VIIB:0.1.3}
If \(A\) and \(B\) are pseudocompact rings, a \emph{homomorphism} from
\(A\) to \(B\) is, by definition, a continuous map compatible with
addition, multiplication, and the identity elements. Such a
homomorphism sends a topologically nilpotent element of \(A\) to a
topologically nilpotent element of \(B\); it therefore maps the radical
\(r(A)\) into the radical \(r(B)\).

